% Generated draft: requires mathematical and source review.
\documentclass[11pt]{article}
\usepackage[margin=25mm]{geometry}
\usepackage{fontspec}
\setmainfont{Latin Modern Roman}
\usepackage{amsmath,amssymb,mathtools,unicode-math}
\setmathfont{Latin Modern Math}
\usepackage{enumitem,needspace,xurl,hyperref}
\setlength{\parindent}{0pt}
\setlength{\parskip}{5pt}
\setlength{\emergencystretch}{3em}
\newcommand{\N}{\mathbb{N}}
\newcommand{\Z}{\mathbb{Z}}
\newcommand{\Q}{\mathbb{Q}}
\newcommand{\R}{\mathbb{R}}
\newcommand{\C}{\mathbb{C}}
\newcommand{\F}{\mathbb{F}}
\newcommand{\A}{\mathbb{A}}
\newcommand{\PP}{\mathbb{P}}
\newcommand{\E}{\mathbb{E}}
\newcommand{\eps}{\varepsilon}
\newcommand{\dd}{\,\mathrm{d}}
\begin{document}
\section*{UnsolvedMath ID: 20000223}
\textbf{Source problem code:} catalogue entry
\subsection{Definitions and mathematical statement}
This record is the contextual introduction to Problem 5, not the separately indexed Problem 5.1. It reports the Bounded Negativity Conjecture in the following source form: for a surface \(X\) that is rational or is over characteristic zero, there is a bound \(n_X\) such that \(C^2\geq n_X\) for every reduced curve \(C\subset X\). Smoothness, projectivity, the field convention and reducible versus irreducible variants are not stated explicitly. The source's comment that nobody knows how to approach it is a dated assessment, not a new literature claim.

For the plane-curve definition, work over an algebraically closed field and let \(C\subset\PP^2\) be reduced and singular, of degree \(d\), with finite nonempty set \(S=\operatorname{Sing}(C)\) of distinct proper singular points. Put \(s=|S|\) and \(m_p=\operatorname{mult}_p(C)\). The broken displayed fraction is reconstructed as
\[
H(C)=h(C)=\frac{d^2-\sum_{p\in S}m_p^2}{s}.
\]
If \(\beta:\widetilde{\PP}^2\to\PP^2\) blows up those distinct proper points once, and \(\widetilde C\) is the proper transform, then \(H(C)=\widetilde C^2/s\). Infinitely near singular points are not included. The denominator excludes smooth curves from this definition.

The attached aid reports an exact characteristic-free numerator gluing formula for arbitrary reduced unions, including tangencies and nonordinary singularities. If \(A,B\) have no common component, all multiplicities in their union can be expressed on
\(T=\operatorname{Sing}(A)\cup\operatorname{Sing}(B)\cup(A\cap B)\), with \(m_A(p)=0\) off \(A\) and similarly for \(B\), by
\[
(\deg A+\deg B)^2-\sum_{p\in T}(m_A(p)+m_B(p))^2.
\]
Its stated transverse specialization is as follows: for singular reduced \(A\) of degree \(a\), with \(s\) singular points, and a smooth degree-\(b\) curve \(B\) avoiding those points and meeting \(A\) transversely at \(ab\) distinct smooth points,
\[
(s+ab)\bigl(H(A\cup B)+2\bigr)=s\bigl(H(A)+2\bigr)+b^2.
\]

After adjoining \(t\geq0\) successive general lines, with all new intersections transverse and distinct and no new triple points, let the resulting curve be \(A_t\). The aid reports
\[
H(A_t)+2=
\frac{s(H(A)+2)+t}{s+at+\binom t2}.
\]
Thus its sign relative to \(-2\) is the sign of \(s(H(A)+2)+t\), giving the reported finite threshold \(t=-s(H(A)+2)\) when that threshold lies in the nonnegative range. As \(t\to\infty\), it reports \(H(A_t)\to-2\), eventually from above. These are supplementary formulas attached to the contextual definition; no further conjectural target is invented for this record.
\subsection{Short English statement}
Reconstruct the plane-curve \(H\)-constant as proper-transform self-intersection divided by the number of distinct proper singular points. This record supplies context for a separately indexed question. The attached aid reports gluing and general-line-padding formulas, including the exact threshold for the padded constant relative to \(-2\).
\subsection{Sources}
\begin{itemize}
\item Local UnsolvedMath registry, permanent ID \(20000223\), exact public code \texttt{catalogue entry}; canonical location \texttt{aim-algebraic-geometry-notes.json}, notes[222].
\item AIM workshop ``Degenerations in algebraic geometry,'' introduction to Problem 5, suggested by Brian Harbourne; supplied source URL \url{https://aimath.org/pastworkshops/degenalggeomproblems.pdf}. The supplied verification note dates the list September 7, 2016 and reports reconstructing the fraction from page 3, PDF index 2. This drafting step did not consult the PDF. The literature field is empty.
\item Attached research attempt 1, ``Exact gluing and transverse-padding formulas for plane-curve H-constants,'' with no complete publication citations.
\end{itemize}
\end{document}
